% GENERATED FILE. Edit canonical lessons, then run npm run generate:books.
% canonical-source-hash: fnv1a64:32516ae527528cff
% canonical-lessons: LA-C513-misericors, LA-C513-misericordia, LA-C513-misereor, LA-C513-pudet, LA-C513-verecundia

\chapter{Pity and shame}
\label{ch:la-pity-and-shame}
\hlchaptermodality{513}

\begin{chapteropening}
\textbf{By the end of this chapter:} \emph{I can talk about pity, compassion and shame.}

Complete the last lesson of chapter 513: I can talk about pity, compassion and shame.
\end{chapteropening}

\section[misericors, misericordis]{misericors, misericordis — compassionate, merciful}

\label{lesson:LA-C513-misericors}



\begin{quote}
Before the new one: say the Latin for praise, then the Latin for to blame.
\end{quote}

\subsection*{You'll want to know: misericors, misericordis}
\textbf{misericors, misericordis} — \textquotedblleft{}compassionate, merciful\textquotedblright{}.

\begin{quote}
Fēmina misericors pauperibus cibum dat.
\end{quote}

\begin{quote}
\emph{The compassionate woman gives food to the poor.}
\end{quote}

One ending: \emph{miser} 'wretched' + \emph{cor} 'heart' — with a heart for the wretched.

\begin{cousinweb}
\textbf{misericors} is built on \emph{miser} \textquotedblleft{}wretched\textquotedblright{} and \emph{cor} \textquotedblleft{}heart\textquotedblright{} (Lewis \& Short; EtymDB).
\end{cousinweb}

\begin{grammarlens}[title={What you've built}]
One more word for telling, explaining and describing at B1.
\end{grammarlens}

\subsection*{Your turn}
\begin{itemize}
\raggedright
  \item \emph{Say it:} \emph{misericors, misericordis}
  \item \emph{Say it:} \emph{misericors, misericordis}, once more
  \item \emph{Say it:} a sentence of your own with \emph{misericors, misericordis}
  \item \emph{Recall:} say the Latin for praise, then the Latin for to blame, then say \emph{misericors, misericordis} again
\end{itemize}

\begin{tcolorbox}[breakable,colback=teal!4,colframe=teal!35!black,title={Before you move on}]
What does misericors, misericordis mean? (\textquotedblleft{}Compassionate, merciful\textquotedblright{}.) Say it once more.
\end{tcolorbox}
\section[misericordia, misericordiae]{misericordia, misericordiae — pity, compassion}

\label{lesson:LA-C513-misericordia}



\begin{quote}
Before the new one: say the Latin for to blame, then the Latin for compassionate.
\end{quote}

\subsection*{You'll want to know: misericordia, misericordiae}
\textbf{misericordia, misericordiae} — \textquotedblleft{}pity, compassion\textquotedblright{}.

\begin{quote}
Misericordiā mōtus, pauperī pecūniam dat.
\end{quote}

\begin{quote}
\emph{Moved by pity, he gives the poor man money.}
\end{quote}

Feminine, first declension, from \emph{misericors}.

\begin{cousinweb}
\textbf{misericordia} is built on \emph{misericors} \textquotedblleft{}compassionate\textquotedblright{}, from \emph{miser} and \emph{cor} \textquotedblleft{}heart\textquotedblright{} (Lewis \& Short; EtymDB).
\end{cousinweb}

\begin{grammarlens}[title={What you've built}]
One more word for telling, explaining and describing at B1.
\end{grammarlens}

\subsection*{Your turn}
\begin{itemize}
\raggedright
  \item \emph{Say it:} \emph{misericordia, misericordiae}
  \item \emph{Say it:} \emph{misericordia, misericordiae}, once more
  \item \emph{Say it:} a sentence of your own with \emph{misericordia, misericordiae}
  \item \emph{Recall:} say the Latin for to blame, then the Latin for compassionate, then say \emph{misericordia, misericordiae} again
\end{itemize}

\begin{tcolorbox}[breakable,colback=teal!4,colframe=teal!35!black,title={Before you move on}]
What does misericordia, misericordiae mean? (\textquotedblleft{}Pity, compassion\textquotedblright{}.) Say it once more.
\end{tcolorbox}
\section[misereor, miserērī]{misereor, miserērī — to pity, feel sorry for}

\label{lesson:LA-C513-misereor}



\begin{quote}
Before the new one: say the Latin for compassionate, then the Latin for pity.
\end{quote}

\subsection*{You'll want to know: misereor, miserērī}
\textbf{misereor, miserērī} — \textquotedblleft{}to pity, feel sorry for\textquotedblright{}.

\begin{quote}
Miserēmur puerōrum, quī nūllōs parentēs habent.
\end{quote}

\begin{quote}
\emph{We pity the children who have no parents.}
\end{quote}

Deponent: \emph{misereor, miserērī, miseritus sum}, with the genitive: \emph{tuī misereor} — I pity you.

\begin{cousinweb}
\textbf{misereor} is built on \emph{miser} \textquotedblleft{}wretched, pitiable\textquotedblright{} (Lewis \& Short). English \emph{misery} comes from \emph{miseria}, from the same adjective.
\end{cousinweb}

\begin{grammarlens}[title={What you've built}]
One more word for telling, explaining and describing at B1.
\end{grammarlens}

\subsection*{Your turn}
\begin{itemize}
\raggedright
  \item \emph{Say it:} \emph{misereor, miserērī}
  \item \emph{Say it:} \emph{misereor, miserērī}, once more
  \item \emph{Say it:} a sentence of your own with \emph{misereor, miserērī}
  \item \emph{Recall:} say the Latin for compassionate, then the Latin for pity, then say \emph{misereor, miserērī} again
\end{itemize}

\begin{tcolorbox}[breakable,colback=teal!4,colframe=teal!35!black,title={Before you move on}]
What does misereor, miserērī mean? (\textquotedblleft{}To pity, feel sorry for\textquotedblright{}.) Say it once more.
\end{tcolorbox}
\section[pudet, pudēre]{pudet, pudēre — it shames, one is ashamed}

\label{lesson:LA-C513-pudet}



\begin{quote}
Before the new one: say the Latin for pity, then the Latin for to pity.
\end{quote}

\subsection*{You'll want to know: pudet, pudēre}
\textbf{pudet, pudēre} — \textquotedblleft{}it shames, one is ashamed\textquotedblright{}.

\begin{quote}
Mē pudet errōris meī.
\end{quote}

\begin{quote}
\emph{I am ashamed of my mistake.}
\end{quote}

Impersonal: \emph{pudet, pudēre, puduit}. The person in the accusative, the cause in the genitive. \emph{Pudor}, already known, is shame.

\begin{cousinweb}
\textbf{pudet} is the impersonal use of \emph{pudeō} \textquotedblleft{}to be ashamed\textquotedblright{}, and \emph{pudor} \textquotedblleft{}shame\textquotedblright{} is built on it (Lewis \& Short). English \emph{impudent} comes from the family.
\end{cousinweb}

\begin{grammarlens}[title={What you've built}]
One more word for telling, explaining and describing at B1.
\end{grammarlens}

\subsection*{Your turn}
\begin{itemize}
\raggedright
  \item \emph{Say it:} \emph{pudet, pudēre}
  \item \emph{Say it:} \emph{pudet, pudēre}, once more
  \item \emph{Say it:} a sentence of your own with \emph{pudet, pudēre}
  \item \emph{Recall:} say the Latin for pity, then the Latin for to pity, then say \emph{pudet, pudēre} again
\end{itemize}

\begin{tcolorbox}[breakable,colback=teal!4,colframe=teal!35!black,title={Before you move on}]
What does pudet, pudēre mean? (\textquotedblleft{}It shames, one is ashamed\textquotedblright{}.) Say it once more.
\end{tcolorbox}
\section[verēcundia, verēcundiae]{verēcundia, verēcundiae — shyness, modesty}

\label{lesson:LA-C513-verecundia}



\begin{quote}
Before the new one: say the Latin for to pity, then the Latin for it shames.
\end{quote}

\subsection*{You'll want to know: verēcundia, verēcundiae}
\textbf{verēcundia, verēcundiae} — \textquotedblleft{}shyness, modesty\textquotedblright{}.

\begin{quote}
Puella prae verēcundiā tacet.
\end{quote}

\begin{quote}
\emph{The girl stays silent out of shyness.}
\end{quote}

Feminine, first declension, from \emph{verēcundus} — shy, modest. A gentle sense of shame.

\begin{cousinweb}
\textbf{verēcundia} is built on \emph{verēcundus} \textquotedblleft{}modest, shy\textquotedblright{}, from \emph{vereor} \textquotedblleft{}to respect, fear\textquotedblright{} (Lewis \& Short; EtymDB).
\end{cousinweb}

\begin{grammarlens}[title={What you've built}]
One more word for telling, explaining and describing at B1.
\end{grammarlens}

\subsection*{Your turn}
\begin{itemize}
\raggedright
  \item \emph{Say it:} \emph{verēcundia, verēcundiae}
  \item \emph{Say it:} \emph{verēcundia, verēcundiae}, once more
  \item \emph{Say it:} a sentence of your own with \emph{verēcundia, verēcundiae}
  \item \emph{Recall:} say the Latin for to pity, then the Latin for it shames, then say \emph{verēcundia, verēcundiae} again
\end{itemize}

\begin{tcolorbox}[breakable,colback=teal!4,colframe=teal!35!black,title={Before you move on}]
What does verēcundia, verēcundiae mean? (\textquotedblleft{}Shyness, modesty\textquotedblright{}.) Say it once more.
\end{tcolorbox}
